\section{Positivity closes the coefficient-to-arithmetic conductor inclusion}
This calculation applies the positivity of the original ES witness to the
entire arithmetic comparison PS1--25 and ISR1--34. The incoming account
\cite[eqs.\ (9)--(23)]{DefiningPrimeReceived} allowed additional residue
collisions. The following bounds determine which of those strata can
occur for actual positive witnesses. The four formal labels and their
eight signs remain the original ones \cite{ES488}; sorting denominator
values in an inequality means an explicitly chosen permutation of the
three denominator labels, never identifying two labels or changing the
quartic. All quantities in the final matrices remain in original order.

\subsection{The complete positive one-divisible bound}
Let $p\equiv1\pmod{12}$, $4/p=1/x+1/y+1/z$, and suppose just $z=pZ$
is divisible by $p$. ISR1--3 proves that $p\nmid xyZ$, that the four
roots are distinct, and that $4Z\equiv1\pmod p$. Choose the permutation
of the two unit denominator labels with $x<y$. Since $p\equiv1\pmod4$,
the least positive residue of $1/4$ is $(3p+1)/4$. Also each denominator
is strictly greater than $p/4$. Consequently
\[
 Z\ge\frac{3p+1}{4},\qquad x\ge\frac{p+3}{4},\qquad
 \frac1y=\frac4p-\frac1x-\frac1{pZ}
 \ge\frac{32}{(p+3)(3p+1)}.
 \tag{PB1}
\]
The last inequality follows by substituting the first two lower bounds
in the two subtracted reciprocals; its numerator is exactly $32p$ before
cancelling the common $p$. Therefore
\[
 0<y-x\le\frac{(p+3)(3p-7)}{32}<p^2,
 \qquad n=v_p(y-x)\in\{0,1\}.
 \tag{PB2}
\]
The strict inequality follows from
$32p^2-(p+3)(3p-7)=29p^2-2p+21>0$.
This argument retains the integer difference itself; a higher valuation
would require the positive difference to be at least $p^2$.
There is no assertion here that every remaining stratum is attained.

\subsection{Every positive two-divisible witness has separated residues}
Now let the only unit denominator be $x$, and write the other two as
$pY,pZ$, choosing the permutation with $Y<Z$. The exact equation is
\[
 \frac{4x-p}{x}=\frac1Y+\frac1Z,\qquad
 Y+Z\equiv4YZ\pmod p,\qquad
 1<Y\le\frac{2x}{4x-p}\le\frac{p+3}{6}<p.
 \tag{PB3}
\]
Here $Y=1$ would make a denominator equal to $p$, excluded in ISR1.
The reciprocal inequality gives the middle bound. The function
$2u/(4u-p)$ has derivative $-2p/(4u-p)^2<0$ for $u>p/4$;
at the least admissible integer $u=(p+3)/4$ it equals $(p+3)/6$.
This proves every displayed inequality.

The residues of $1,Y,Z$ are pairwise distinct. First $Y\not\equiv1$
because $1<Y<p$. If $Z\equiv1$, the original congruence would give
$3Y\equiv1$. Its least positive solution, since $p\equiv1\pmod3$,
is $(2p+1)/3$, strictly above $(p+3)/6$. If $Z\equiv Y$, it would
give $2Y\equiv1$ since $Y$ is a unit. Its least positive solution
is $(p+1)/2$, also strictly above $(p+3)/6$. These exhaust all possible
pairwise residue coincidences. Thus, in the original labelled cluster,
\[
 v_p(p-pY)=v_p(p-pZ)=v_p(pY-pZ)=1.
 \tag{PB4}
\]
In particular the closer-pair cases $n\ge2$ in the general ISR17
calculation never occur for a positive witness. Positivity removes them;
the validity of the more general matrix calculation is unchanged.

\subsection{The full ideal inclusion and its least coefficient power}
Continue with the original integral chart $p\nmid S$, so $A=-1/S$
is a unit. No coordinate or coefficient is changed to impose this
condition; ISR18--23 treats the different fractional chart explicitly.
Let $I=(E:O)$ be the complete algebra conductor PS6--12 and let $C_0$
be the original quartic coefficient. Up to the stated permutation,
the only possible lists are
\[
\begin{array}{c|c|c|c|c}
 \text{divisible denominators}&n&(b_i)&(3m_i+\epsilon_i)&
              \text{Smith exponents of }O/E\text{ and }E/I\\\hline
 1&0&(1,0,0,1)&(1,0,0,1)&(0,0,0,0,0,0,1,1)\\
 1&1&(1,1,1,1)&(1,1,1,1)&(0,0,0,0,1,1,1,1)\\
 2&1&(2,0,2,2)&(3,0,3,3)&(0,0,0,1,1,2,2,3).
\end{array}\tag{PB5}
\]
For the first two rows the label order is $(p,x,y,z)$; for the last
row it is $(p,x,pY,pZ)$. The complete full units remain
$u_i=p^{-b_i}Af'(r_i)$ and the ideal is still
$I=\prod_i p^{3m_i+\epsilon_i}u_iO_i$, not just its list of valuations.
PB5 follows by substituting PB2 and PB4 into the fully proved ISR16--17.
The two-divisible quadratic factors are all unramified or split,
because their derivative valuations are even.

The original coefficient is $C_0=-5Axyz$, so its valuation is respectively
one or two. For every integer $j\ge0$, PS9 gives the exact equivalence
\[
 C_0^jE\subset I\quad\Longleftrightarrow\quad
 jv_p(C_0)\ge\max_i(3m_i+\epsilon_i).
 \tag{PB6}
\]
Indeed $I$ is an $E$-ideal containing $C_0^jE$ exactly when it contains
the element $C_0^j$, and its branch ideals give the displayed criterion.
Combining PB5--6 proves, for every positive witness in this integral chart,
\[
 \boxed{\Pi_P(C_0^2E_N)=C_0^2E\subset I.}\qquad
 \min\{j\ge0:C_0^jE\subset I\}=
 \begin{cases}1,&\text{one divisible denominator},\\
 2,&\text{two divisible denominators}.
 \end{cases}\tag{PB7}
\]
The equality is the full surjective specialization PS14; it is not an
identification of the two different conductor ideals. The minima hold
because the largest branch exponent in PB5 is one or three, respectively.
Thus the coefficient conductor inclusion, formerly restricted by an
unevaluated gap condition, now holds for the whole actual positive
witness class in its integral chart.

All matrices realizing this inclusion in the original conductor remain
explicit: $F_UJ_I,F_WJ_I$ and $F_U\Pi_P(C_0^2I),F_W\Pi_P(C_0^2I)$,
with the unchanged original receivers PS23--25. For a direct eight-column
factorization put $J=J_I$ and define
\[
 Z_j=J^{-1}C_0^jI_8
 =\diag\!\left(C_0^j A^{-1}\diag(p^{-m_i})Q^{-1},
                   C_0^j A^{-1}Q^{-1}\right),\qquad
 JZ_j=C_0^jI_8.
 \tag{PB8}
\]
Every entry is integral for the values of $j$ in PB7 by the proved
lattice containment, and the displayed inverse specifies the full
matrix without losing any direction. Multiplication by the original
receivers gives $(F_bJ)Z_j=C_0^jF_b$ for $b=U,W$, and the two
identities intertwine by $T_A\oplus T_A$ exactly as in PS23.
The norms and quotient minima remain those in PS24--25. In particular
PB7 closes an arithmetic ideal comparison; it does not replace the
original linear conductor by a multiplication operator or force its
first nonzero moment to vanish.
